CCD cameras on space probes are very sensitive and exposed to space weather
conditions. Intense radiation passing through the CCD produces electron-hole
pairs in the silicon of the CCD chip. The rate at which these pairs are
produced is an important parameter when operating cameras on board spacecraft
and can be calculated for radiation of any given energy.

A high energy particle or photon of radiation passing through the CCD will
deposit some energy in the chip with each electron-hole pair it creates. The
`stopping power' of silicon for a given type of particle can be measured as
the energy per areal density ($\mathrm{areal\ density} = \mathrm{mass\ per\
unit\ area}$) that the silicon `takes away' from the travelling particle.

\ioaafig{2021_TH_Q13-f1.pdf}

NASA's Lucy mission will be the first to study the Trojan asteroids and will
revolutionize our understanding of the formation of the Solar System. One of
the instruments on board is L'LORRI (Lucy LOng Range Reconnaissance Imager),
which contains a sensitive CCD in order to produce detailed images of the
Trojan asteroids. Unfortunately, the radiation around Jupiter is very intense
and it can generate a lot of `noise' in the pixels of the CCD.

Let us assume that an average charged particle trapped in Jupiter's magnetic
field has an energy of 15\,MeV and that the flux of such particles in this
region is equivalent to about 600 electrons $\mathrm{s}^{-1}\,
\mathrm{cm}^{-2}$. Also assume that for each electron-hole pair which a
particle passing through a pixel creates, it deposits exactly the excitation
energy of the pair in that pixel. After the pixel crosses a threshold number
of electron-hole pairs it is `excited' and no more pairs can be produced in
that pixel. Any remaining energy in the particle is passed to the next pixel
(and so on).

Using the data given below for the CCD chip in the L'LORRI camera, answer the
following questions:

\begin{description}
\item[(13.1)] How many pixels will be excited by one such particle of
  radiation passing through the CCD when the spacecraft is near Jupiter's
  orbit? \hfill[10.0\,pt]
\item[(13.2)] Given the radiation flux near Jupiter, what percentage of the
  total number of pixels in an image will be excited? \hfill[5.0\,pt]
\end{description}

CCD Data:
\begin{itemize}
\item Exposure time of an image $= 30$\,ms
\item Pixels on the CCD $= 1024 \times 1024$
\item CCD Area $= 13\,\mathrm{mm} \times 13\,\mathrm{mm}$
\item CCD chip thickness $= 0.06$\,cm
\item Density of silicon, $\rho = 2.34\,\mathrm{g\,cm}^{-3}$
\item Excitation energy of single pair $= 2.36$\,eV
\item Excitation threshold of a single pixel $= 250$ pairs
\item `Stopping power' of silicon for a 15\,MeV electron
  $= 3.012\,\mathrm{MeV\,g}^{-1}\,\mathrm{cm}^{2}$
\end{itemize}
